\topic{25}{Dyfrakcja światła i zasada Huygensa--Fresnela}{
Zasada Huygensa--Fresnela traktuje punkty czoła fali
jako źródła fal wtórnych, których superpozycja
wyznacza pole dalej. Dyfrakcja staje się widoczna,
gdy rozmiar przeszkody jest porównywalny
z długością fali. W obszarze Fraunhofera
obserwujemy rozkład w dalekim polu lub
w płaszczyźnie ogniskowej soczewki.

Dla szczeliny szerokości $a$ minima natężenia
spełniają $a\sin\theta=m\lambda$ dla
całkowitego $m\ne0$.
Centralne maksimum jest szersze niż boczne.
Dla siatki o stałej $d$ główne maksima
spełniają $d\sin\theta=m\lambda$.
Skończona liczba szczelin określa szerokość
prążków, a ich szerokość własna kształtuje
obwiednię. Równania zakładają spójne
oświetlenie i odpowiednią geometrię.}
\topic{26}{Optyka geometryczna i podstawowe przyrządy}{
Optyka geometryczna jest granicą falową,
w której długość fali jest mała względem
rozmiarów układu. Promień wskazuje kierunek
przenoszenia energii, a soczewka cienka
spełnia $1/f=1/p+1/q$ przy konsekwentnej
konwencji znaków. Powiększenie liniowe
wynosi $M=-q/p$.

Lupa zwiększa kąt widzenia małego przedmiotu;
jej przybliżone powiększenie kątowe zależy
od ogniskowej i punktu bliskiego oka.
Luneta łączy obiektyw tworzący obraz
dalekiego obiektu z okularem, który ten obraz
ogląda. Mikroskop najpierw tworzy powiększony
obraz pośredni obiektywem, a następnie
powiększa kątowo przez okular.
Rozdzielczości nie można zwiększać dowolnie:
ogranicza ją dyfrakcja i apertura układu.}
\newpage
\topic{27}{Odbicie, załamanie, całkowite odbicie i dyspersja}{
Na granicy ośrodków kąt odbicia równa się
kątowi padania. Prawo Snelliusa
$n_1\sin\theta_1=n_2\sin\theta_2$
wynika z ciągłości fazy wzdłuż granicy.
Przy przejściu z większego $n$ do mniejszego
całkowite wewnętrzne odbicie zachodzi
dla $\theta_1>\theta_c$, gdzie
$\sin\theta_c=n_2/n_1$.

Współczynnik załamania zależy od częstotliwości.
Dyspersja normalna oznacza zwykle wzrost $n$
z częstością w obszarze dalekim od rezonansu;
w pobliżu pasma absorpcyjnego może wystąpić
dyspersja anomalna. Szybkość fazowa $c/n$
nie musi równać się grupowej. W światłowodzie
całkowite odbicie pomaga prowadzić światło,
ale rzeczywiste straty zależą też od absorpcji
i rozpraszania.}
\topic{28}{Polaryzacja, prawo Malusa i dwójłomność}{
Polaryzacja opisuje kierunek drgań wektora
pola elektrycznego poprzecznej fali.
Idealny polaryzator przepuszcza składową
wzdłuż swojej osi; dla światła liniowo
spolaryzowanego prawo Malusa daje
$I=I_0\cos^2\theta$. Dla światła
niespolaryzowanego pierwszy idealny
polaryzator przepuszcza średnio połowę
natężenia.

Przy padaniu pod kątem Brewstera światło
odbite może być liniowo spolaryzowane;
$\tan\theta_B=n_2/n_1$ dla ośrodków
niemagnetycznych. Dwójłomny kryształ
ma różne współczynniki załamania dla
różnych polaryzacji. Dwójłomność może
też być wymuszona naprężeniem lub polem.
Aktywność optyczna obraca płaszczyznę
polaryzacji i wiąże się m.in. z chiralnością
ośrodka.}
\newpage
\topic{29}{Naturalna szerokość linii i stany metatrwałe}{
Stan wzbudzony o skończonym czasie życia
$\tau$ ma niezerową szerokość energetyczną.
Dla zaniku wykładniczego profil linii
jest lorentzowski, a rząd wielkości
szerokości energii to $\Delta E\sim\hbar/\tau$;
precyzyjny współczynnik zależy od przyjętej
definicji szerokości. Krótki czas życia
oznacza szerszą linię.

W doświadczeniu obserwowana szerokość może
być większa z powodu efektu Dopplera,
zderzeń lub aparatury. Stan metatrwały
ma wyjątkowo długi czas życia, ponieważ
przejście do niższego stanu jest słabo
dozwolone przez reguły wyboru lub
wymaga wyższego rzędu procesu.
Takie stany umożliwiają gromadzenie obsadzenia
w laserach. Trzeba odróżnić naturalną
szerokość pojedynczego przejścia od
poszerzenia całego badanego widma.}
\topic{30}{Doświadczenie Rutherforda i model atomu}{
Rutherford badał rozpraszanie cząstek alfa
na cienkiej folii. Większość przelatywała
z niewielkim odchyleniem, ale mała część
ulegała dużym kątom rozproszenia.
Model dodatniego ładunku rozłożonego
równomiernie w całej objętości atomu
nie wyjaśniał łatwo tych rzadkich
silnych odchyleń.

Wniosek był taki, że dodatni ładunek
i niemal cała masa są skupione w bardzo
małym jądrze. Klasyczny wzór Rutherforda
na różniczkowy przekrój czynny ma zależność
$d\sigma/d\Omega\propto\sin^{-4}(\theta/2)$
dla odpychania kulombowskiego cząstek
punktowych. Model jądrowy rozwiązywał
problem rozpraszania, lecz klasyczna
elektrodynamika nie wyjaśniała stabilności
orbit elektronowych; wymagała ona
opisu kwantowego.}
\newpage
\topic{31}{Widmo promieniowania rentgenowskiego}{
Promieniowanie rentgenowskie powstaje
gdy szybkie elektrony hamują w materiale
anody albo wybijają elektrony z wewnętrznych
powłok atomu. Hamowanie daje widmo ciągłe.
Jego graniczna długość fali spełnia
$\lambda_{\min}=hc/(eU)$, gdy elektron
przyspieszono napięciem $U$ i całą
energię oddał jednemu fotonowi.

Promieniowanie charakterystyczne tworzy
linie związane z przejściami między
powłokami po powstaniu wakatu.
Ich energie zależą od pierwiastka anody,
podczas gdy granica widma ciągłego
zależy od napięcia. Intensywność
i kształt widma zależą również od
filtracji i absorpcji w materiale.
Przy opisie widma trzeba rozróżnić
proces wytwarzania fotonów od ich
późniejszego osłabienia w ośrodku.}
\topic{32}{Wzbudzanie i deekscytacja jąder oraz atomów}{
Wzbudzone jądro może oddać energię
jako kwant gamma albo przez konwersję
wewnętrzną: energia przejścia trafia
bezpośrednio do elektronu atomowego.
Po powstaniu wakatu w powłoce atom
może wyemitować promieniowanie X
lub elektron Augera. Wybór kanału
zależy od energii i struktury atomu.

Efekt Mössbauera polega na niemal
bezodrzutowej emisji i absorpcji
kwantów gamma przez jądra osadzone
w sieci krystalicznej. Cała sieć
przejmuje pęd, więc energia odrzutu
pojedynczego jądra staje się znikoma.
Umożliwia to bardzo wąskie rezonanse
i badanie subtelnych przesunięć
energetycznych. Należy odróżniać
konwersję wewnętrzną od rozpadu beta:
w pierwszym przypadku jądro
nie zmienia liczby protonów.}
\newpage
\topic{33}{Jądro atomowe: odkrycie, budowa i własności}{
Doświadczenie Rutherforda wskazało
istnienie małego, dodatnio naładowanego
jądra. Składa się ono z protonów
i neutronów, czyli nukleonów.
Liczba protonów $Z$ określa pierwiastek,
a liczba masowa $A=Z+N$ sumę nukleonów.
Izotopy mają to samo $Z$, lecz różne $N$.

Typowy promień jądra przybliża wzór
$R=r_0A^{1/3}$, gdzie $r_0$ jest rzędu
femtometra. Oznacza to prawie stałą gęstość
materii jądrowej. Masa związanego jądra
jest mniejsza od sumy mas swobodnych
nukleonów; defekt masy odpowiada
energii wiązania $E_b=\Delta mc^2$.
Jądra mają spin, moment magnetyczny
i dyskretne stany wzbudzone.
Stabilność wynika z równowagi sił
jądrowych i odpychania Coulomba.}
\topic{34}{Model kroplowy i powłokowy jądra}{
Model kroplowy traktuje jądro
jak naładowaną kroplę o niemal stałej
gęstości. Półempiryczny wzór masowy
zawiera składniki objętościowy,
powierzchniowy, Coulomba, asymetrii
i parowania. Pomaga wyjaśnić
ogólne trendy energii wiązania
oraz możliwość rozszczepienia
ciężkich jąder.

Model powłokowy opisuje nukleony
poruszające się w średnim potencjale
z uwzględnieniem zakazu Pauliego.
Silne sprzężenie spin--orbita
prowadzi do liczb magicznych
i zwiększonej stabilności
zamkniętych powłok. Model
wyjaśnia spiny i momenty wielu
jąder, których sama kropla
nie opisuje. Oba modele
są przybliżeniami o różnych
obszarach skuteczności;
jądro może wykazywać zarazem
cechy kolektywne i indywidualne.}
\newpage
\topic{35}{Spontaniczne przemiany jądrowe}{
Rozpad alfa emituje jądro helu i jest
istotny dla ciężkich nuklidów;
bariera Coulomba jest pokonywana
przez tunelowanie. W rozpadzie beta
oddziaływanie słabe zmienia neutron
w proton lub odwrotnie, z emisją
leptonów zapewniającą zachowanie
ładunku i liczby leptonowej.
Rozpad gamma usuwa energię wzbudzenia
bez zmiany $A$ i $Z$.

Prawo zaniku pojedynczego kanału
ma postać $N(t)=N_0e^{-\lambda t}$,
a czas połowicznego zaniku
$T_{1/2}=\ln2/\lambda$.
Różne kanały mają częściowe
stałe zaniku sumujące się
do całkowitej. Energia wydzielona
$Q$ wynika z różnicy mas
początkowych i końcowych.
Przemiana jest spontaniczna,
gdy bilans energii pozwala
na stan końcowy, ale szybkość
może być bardzo mała.}
\topic{36}{Oddziaływanie promieniowania jądrowego z materią}{
Cząstki naładowane tracą energię
głównie przez jonizację i wzbudzanie
atomów. Ciężkie cząstki alfa
mają krótki zasięg i gęsto
jonizują; elektrony beta ulegają
też silnemu rozpraszaniu
i mogą wytwarzać promieniowanie
hamowania. Fotony gamma oddziałują
m.in. przez efekt fotoelektryczny,
rozpraszanie Comptona i tworzenie
par przy odpowiedniej energii.

Neutrony są elektrycznie obojętne,
więc oddziałują z jądrami
przez rozpraszanie i reakcje
jądrowe. Osłabienie wąskiej
wiązki fotonów bywa opisane
$I(x)=I_0e^{-\mu x}$,
ale współczynnik zależy od
energii i materiału. Zasięg
cząstki i współczynnik osłabienia
fotonów to różne wielkości;
przy ochronie radiologicznej
trzeba uwzględnić również
promieniowanie wtórne.}
\newpage
\topic{37}{Rozszczepienie jąder i energetyka jądrowa}{
Ciężkie jądro może rozpaść się
na dwa lżejsze fragmenty.
W rozszczepieniu indukowanym
neutronem powstają zwykle
dodatkowe neutrony, co pozwala
na reakcję łańcuchową.
Energia wydzielona wynika
z większej energii wiązania
na nukleon w produktach:
$Q=(m_{\rm przed}-m_{\rm po})c^2$.

Reaktor utrzymuje średnio
jedną następną reakcję
z każdego rozszczepienia
w stanie krytycznym.
Moderator spowalnia neutrony
w reaktorach termicznych,
a pręty sterujące zmieniają
ich pochłanianie. Ciepło
z reakcji jest przenoszone
do obiegu energetycznego
i zamieniane na pracę
przez turbinę. Bilans
techniczny obejmuje sprawność
cieplną, chłodzenie,
odpady promieniotwórcze
i bezpieczeństwo układu.}
\newpage
